% =========================================================
%  frontmatter/petunjuk-penggunaan.tex
%  Cara menggunakan buku ajar
% =========================================================

\chapter*{Petunjuk Penggunaan Buku}
\addcontentsline{toc}{chapter}{Petunjuk Penggunaan Buku}
\thispagestyle{plain}

Bagian ini menjelaskan bagaimana buku ini disusun dan bagaimana cara
menggunakannya agar kegiatan belajar berjalan efektif.

\section*{Hubungan Buku dengan Rencana Pembelajaran Semester}

Buku ini bukan dokumen yang berdiri sendiri. Isinya mengikuti Rencana
Pembelajaran Semester (RPS) mata kuliah \NamaMK\ (\KodeMK) yang berlaku di
Program Studi \NamaProdi. Setiap bab inti pada buku ini setara dengan satu
materi pembelajaran pada RPS, dan satu materi tersebut menopang satu atau
beberapa Sub-CPMK. Keterkaitan lengkap antara CPL, CPMK, Sub-CPMK, bab, dan
minggu perkuliahan disajikan pada bagian \istilah{Capaian Pembelajaran dan
Peta Kompetensi} yang menyusul sebelum Bab I.

Karena itu, urutan bab sebaiknya diikuti sebagaimana adanya. Materi
Python/Pygame pada Bab III sampai Bab XI menjadi prasyarat bagi materi
Unity pada Bab XII sampai Bab XVI, dan Bab XVII baru bermakna setelah
seluruh bab sebelumnya dikerjakan.

\section*{Struktur Setiap Bab}

Bab inti pada buku ini selalu memuat bagian berikut dengan urutan yang
tetap, sehingga pembaca dapat menemukan informasi tertentu dengan cepat.

\begin{center}
\begin{tabularx}{\textwidth}{@{}L{4.3cm} Y@{}}
\toprule
\headrow \textbf{Bagian} & \textbf{Fungsi} \\
\midrule
Sub-CPMK Bab Ini &
Rumusan kompetensi yang harus dicapai beserta indikator dan tingkat
\istilah{Taksonomi Bloom}-nya. Bagian ini menjadi dasar penilaian bab. \\
\addlinespace
Tujuan Pembelajaran &
Uraian singkat apa yang mampu dilakukan pembaca setelah menuntaskan bab. \\
\addlinespace
Pendahuluan &
Pengait dengan bab sebelumnya dan peta isi bab. \\
\addlinespace
Materi Pokok &
Pembahasan utama, terbagi menjadi beberapa subbab bernomor. \\
\addlinespace
Contoh &
Contoh kerja yang dapat diikuti langkah demi langkah. \\
\addlinespace
Praktikum &
Lembar kerja yang harus dikerjakan pada komputer, memakai kode program
yang tersedia pada lampiran daring dan repositori mata kuliah. \\
\addlinespace
Latihan &
Soal latihan dan butir refleksi untuk menguji pemahaman sendiri. \\
\addlinespace
Rangkuman &
Ringkasan butir-butir penting bab. \\
\addlinespace
Checklist Pencapaian Kompetensi &
Daftar indikator yang dapat dicentang pembaca untuk menilai kesiapan
sebelum asesmen. \\
\addlinespace
Asesmen dan Instrumen Penilaian &
Bentuk asesmen bab beserta rubrik atau kunci penilaiannya. \\
\bottomrule
\end{tabularx}
\end{center}

\section*{Kotak Informasi}

Sepanjang buku ini dipakai beberapa kotak dengan makna yang tetap.

\begin{itemize}
  \item \textbf{Tujuan Pembelajaran} dan \textbf{Sub-CPMK Bab Ini}:
        kompetensi dan indikator yang menjadi sasaran bab.
  \item \textbf{Contoh}: contoh kerja yang sebaiknya diketik ulang, bukan
        hanya dibaca.
  \item \textbf{Catatan}: penjelasan tambahan atau pengait dengan konsep
        lain.
  \item \textbf{Perhatian}: kesalahan yang sering terjadi dan perlu
        dihindari.
  \item \textbf{Tip Praktik}: saran praktis yang mempercepat pekerjaan.
  \item \textbf{Praktikum}: lembar kerja yang harus diselesaikan.
  \item \textbf{Checklist Pencapaian Kompetensi}: daftar periksa mandiri.
  \item \textbf{Refleksi}: pertanyaan untuk meninjau kembali cara belajar.
\end{itemize}

\section*{Konvensi Penulisan}

\begin{itemize}
  \item Istilah asing dicetak miring, misalnya \istilah{game loop} dan
        \istilah{collision}.
  \item Nama berkas, perintah, nama fungsi, dan potongan kode di dalam teks
        dicetak dengan huruf mesin tik, misalnya \code{pygame.time.Clock}.
  \item Kode program yang panjang disajikan sebagai daftar kode bernomor
        dan disimpan sebagai berkas terpisah pada folder \code{code/} proyek
        buku, sehingga dapat dijalankan langsung di komputer.
  \item Rujukan pustaka ditandai nomor dalam kurung siku, misalnya
        \mbox{[R3]}, dan dapat dicari pada Daftar Pustaka di bagian akhir
        buku. Nomor rujukan ini sama dengan nomor rujukan pada RPS mata
        kuliah, sehingga penelusuran antara kedua dokumen menjadi mudah.
\end{itemize}

\section*{Cara Belajar yang Disarankan}

\begin{langkah}
  \item Bacalah bagian \istilah{Sub-CPMK Bab Ini} dan \istilah{Tujuan
        Pembelajaran} lebih dahulu untuk mengetahui sasaran bab.
  \item Bacalah materi pokok secara aktif: ketik ulang setiap potongan kode
        yang penting, lalu ubah nilainya untuk melihat pengaruhnya.
  \item Kerjakan praktikum sampai program berjalan. Jika program gagal,
        bandingkan kembali dengan daftar kode pada bab dan periksa pesan
        kesalahan yang muncul.
  \item Kerjakan latihan tanpa melihat jawaban terlebih dahulu, lalu
        gunakan checklist kompetensi untuk menilai diri sendiri.
  \item Simpan setiap hasil praktikum pada repositori pribadi. Artefak
        inilah yang menjadi bukti capaian pembelajaran pada penilaian
        portofolio.
  \item Jawab butir refleksi pada akhir bab secara jujur. Bagian yang terasa
        sulit menunjukkan apa yang perlu diulang.
\end{langkah}

\begin{catatabox}
Buku ini menuntut pekerjaan pada komputer. Pastikan lingkungan
pengembangan sudah siap sebelum memulai Bab III. Langkah pemasangan
Python, Pygame, dan editor diuraikan pada bab tersebut, dan pemasangan
Unity diuraikan pada Bab XII.
\end{catatabox}

\clearpage
